\section{Introduction}
Blackjack is a popular casino game where players compete individually against the dealer (the house).\\
In this game, the house generally maintains a slight advantage over the players and for this reason players often employ 
strategies, such as card counting, to mitigate this disadvantage. \\
Card counting involves monitoring the sequence of high and low-value cards already dealt to estimate the 
composition of the remaining deck. 
By maintaining a "running count," a player can identify advantageous moments to adjust their betting patterns. 
\\ \newline
This project proposes the development of an automated card counting detection system. \\
The primary goal is to build a video analysis pipeline capable of interpreting the flow of cards dealt during a game.
Through this process, the system is able to monitor the game status and identify situations favorable to the players, alerting
to focus on players' decisions and betting variations.
\\ \newline
The project focuses on the implementation of the Hi-Lo counting method, where each card is assigned a specific value 
to update the total count:
\begin{itemize}
    \item Low-value cards (2-6): assigned a value of +1;
    \item Neutral cards (7-9): assigned a value of 0;
    \item High-value cards (10-Ace): assigned a value of -1.
\end{itemize}
To facilitate intuitive monitoring, the system overlays color-coded bounding boxes around detected cards 
(Green for +1, Blue for 0, and Red for -1), allowing for real-time interpretation of the counting logic.
\\ \newline
This project aims to show how computer vision can automate blackjack tracking. 
By using real-time detection and Hi-Lo logic, the system monitors the deck and visualizes the count, demonstrating 
the potential for automated monitoring in gaming environments.

\subsection{General Information, Dependencies and Running the project}
The Github repo documenting the project is provided at the following link: \cite{github_repo}. It is required to satisfy the following dependencies:
\begin{itemize}
    \item CMake version: 4.0.0+
    \item OpenCV version: 4+
    \item ONNX Runtime version: 1.21.0. The binaries for the CPU version are already bundled inside the project inside the external/ folder and CMAKE is already configured to find the binaries either in this folder or in the system's directories. The GPU version can be downloaded from the most recent Github Release and it's compatible with CUDA 12.X versions. If both packages are present in the external/ folder, the GPU version will be selected first.
\end{itemize}
Please follow the instructions for running the project in the README.md file.